\section{\rev{제안 방법과 구현}}
\label{sec:method}

\begin{revision}

\subsection{자동 텍스트 변환과 작성 프레임}
\label{sec:extraction}
기존 자동 변환은 LLM 호출이나 수작업 정답 입력을 쓰지 않는 정규식 기반 영어 파서다. \texttt{stop/halt/hold}, \texttt{yield/decelerate/slow down}, \texttt{accelerate/proceed/resume/move forward}, \texttt{maintain/keep speed}를 행동 종류로 묶는다. 명령형ㆍ조동사ㆍ주어 문맥을 확인하여 단순 언급을 행동으로 해석하는 것을 제한한다. \texttt{then/and then}으로 단계 순서를 나누되 \texttt{until then}은 단계 경계로 자르지 않는다. \texttt{until} 뒤 구절은 해제 조건, \texttt{once/after} 뒤 구절은 계기로 보존하고 조건 구절 내부의 행동 단어는 제외한다.

정지ㆍ양보 대상이 없으면 \texttt{UNKNOWN\_MISSING\_TARGET}, 명확한 순서 없는 복수 행동이면 \texttt{UNKNOWN\_AMBIGUOUS\_ORDER}로 둔다. 명시적 단계가 각각 단일 행동일 때만 다음 행동을 연결한다. 부정 표현을 참ㆍ거짓 관측 증거로 읽거나 동의 표현을 범용 의미로 해석하는 기능은 없다. \texttt{release\_known}과 \texttt{satisfaction\_known}은 모두 false로 시작한다. 문구를 찾은 것과 해당 조건이 실제로 충족되었음을 안 것은 다르다.

작성 프레임은 설계자가 원문과 동시에 만든다. \texttt{action}, \texttt{obligation\{id,target,condition\_id\}}, \texttt{evidence\{condition\_id,value\}}, \texttt{precondition}과 \texttt{source\_spans}를 가진다. True/False는 작성 문장이 명시한 조건 상태이며 사람의 독립 센서 정답이 아니다. 기존 모델에 입력할 때 조건 A를 단일 해제 채널에 투영하고 B/OTHER를 A로 바꾸지 않는다. 두 번째 의무 C의 별도 해제 정보를 담지 못한 손실도 기록한다. 조건별 모델은 프레임의 각 의무를 직접 인코딩한다.

\begin{table*}[t]
\centering
\color{revisionblue}
\bitablecaption{가상 보행자 사례의 원문--작성 프레임--모델 상태--판정 연결(중간 사건 0개).}{End-to-end authored pedestrian example with zero intervening events.}
\label{tab:worked-example}
\scriptsize
\begin{tabularx}{\textwidth}{L{.20\textwidth}L{.32\textwidth}L{.25\textwidth}Y}
\toprule
사건/변형 & 실제 작성 원문 구절 & 프레임과 상태 & 출력 \\
\midrule
공통 e0 & Stop until the pedestrian has passed. The required mirror check has been completed. & STOP; 의무 A, 대상 pedestrian; 해제조건 A; precondition=T; active[A]=true, known[A]=? & 의무 시작; 문제 없음 \\
미해제 e1 & The pedestrian is still crossing. The traffic signal is green. Proceed. & GO; evidence A=F, B=T; active[A]=true & 모순; e1, A, UNRELEASED\_OBLIGATION \\
증거 누락 e1 & The traffic signal is green. Proceed. & GO; evidence B=T; known[A]=? & 근거 부족; e1, A, RELEASE\_EVIDENCE\_MISSING \\
정상 해제 e1 & The pedestrian has passed. The traffic signal is green. Proceed. & GO; evidence A=T, B=T; active[A]=false & 해당 의무와 양립 \\
\bottomrule
\end{tabularx}
\end{table*}

표~\ref{tab:worked-example}의 원문은 저장된 작성 사례에서 가져왔다. 미해제 e1의 \texttt{Proceed}는 문자 위치 [63,70), A 증거는 [0,33), B 증거는 [34,62)에 연결된다. 사건 ID는 \texttt{pedestrian-held\_local\_conflict-gap0:e1}이다. 누락ㆍ정상 변형도 각각의 case ID에 e1을 연결한다. 실제 입력은 UTF-8 구절의 문자열 인덱스이며 이 예의 원문은 ASCII다. 자동 파서에 같은 문장을 넣는 경로는 증거를 확정하지 않으므로 작성 프레임의 성공과 혼합하지 않는다.

\subsection{기존 검사기와 보완 모델의 차이}
\label{sec:models}
기존 Python은 활성 의무 한 개, 행동 4종, 상태 6종과 누적 오류ㆍ미상 사유를 관리한다. 단계마다 해석 미상, 알려진 증거 갱신, 오류 규칙, 상태 전이, 누적 출력 순으로 처리한다. 복수 의무는 \texttt{OVERLAPPING\_OBLIGATION}, 대기 중 해제 미상도 누적 미상으로 남을 수 있다. 기존 UPPAAL은 같은 전이 규약과 해시를 사용한다. 이 두 경로는 동결하여 최초 실패를 보존했다.

보완 모델은 실패 분석 뒤 별도 개발했다. 각 의무의 \texttt{active[K]}, 조건별 \texttt{known[K]}(-1/0/1), 선행조건, 사건 인덱스, \texttt{bad/missing} 누적 플래그와 최초 위치ㆍ의무를 둔다. 입력의 의무 등록 순서가 배열 순서이며 선행조건 위반을 먼저, 의무를 등록 순서대로 판정한다. 같은 사건의 여러 위반을 개별 사유로 누적하되 최초 의무 출력은 이 우선순위로 하나를 선택한다. 모든 문제의 전체 순위나 모든 관련 의무 목록을 반환하는 기능은 아니다.

\begin{table}[t]
\centering
\color{revisionblue}
\bitablecaption{보완 모델의 조건별 상태 갱신과 기록 순서.}{Update and recording order in the condition-indexed model.}
\label{tab:transition}
\scriptsize
\begin{tabularx}{\columnwidth}{L{.27\columnwidth}Y}
\toprule
단계 & 실행과 기록 \\
\midrule
1. 등록 & 현재 사건의 새 의무를 active에 넣음 \\
2. 증거 & 알려진 선행조건과 조건별 증거만 갱신; 미상은 이전 값 보존 \\
3. 해제 & 각 의무의 known=1이면 해당 active만 false \\
4. 진행 & GO일 때 precondition=0이면 bad; 이후 등록 순서로 active/known=0이면 bad, -1이면 missing \\
5. 최초 기록 & 최초 bad/missing의 사건ㆍ의무 인덱스를 한 번만 기록; 이후 플래그 누적 \\
6. 종료 & 사건 배열 끝에서 Done; 미래 해제의 부재만으로 오류를 추가하지 않음 \\
\bottomrule
\end{tabularx}
\end{table}

\begin{lstlisting}[caption={보완 모델의 핵심 실행 순서(실제 XML의 step 함수 요약).},label={lst:condition-step}]
activate(creates[idx]);
update_known_precondition_if_present();
for j in obligation_registration_order:
    if observation(idx,j) >= 0: known[j] = observation(idx,j)
    if active[j] and known[j] == 1: active[j] = false
if actions[idx] == GO:
    record_bad_if(precondition == 0, PRECONDITION)
    for j in obligation_registration_order:
        record_bad_if(active[j] and known[j] == 0, j)
        record_missing_if(active[j] and known[j] == -1, j)
idx = idx + 1
\end{lstlisting}

WAIT는 작성 자료에서 유지ㆍ대기 단계를 나타내는 기호다. 실제 속도나 이동을 적분하지 않으므로 \texttt{Maintain speed}라는 문구를 물리적으로 안전한 정지라고 보장하지 않는다. 이 단계는 명시된 GO 전환 위반을 검사하기 위한 중간 사건이다.

\subsection{UPPAAL 상태ㆍ질의와 추적}
\label{sec:queries}
보완 모델은 committed 위치 Run에서 \texttt{idx<N}일 때 \texttt{step(),idx++}로 자기 전이하고, \texttt{idx==N}이면 Done으로 이동한다. committed 처리로 한 사건의 갱신을 끝내고 다음 사건을 읽는다. 정량 시간 시계는 없다. 다음 다섯 질의를 실제 verifyta 5.0.0에 실행했다. 경로ㆍ상태 양화사의 의미는 공식 문서에 따른다 \cite{uppaalqueries}.

\begin{lstlisting}[caption={보완 모델의 실제 queries.q.},label={lst:queries}]
A[] not bad
A[] not missing
A<> Observer.Done
E<> Observer.Done && bad_hold
E<> Observer.Done && bad_pre
\end{lstlisting}

첫 질의의 실패는 모순, 둘째 질의만 실패하면 근거 부족이다. 셋째는 전체 배열의 처리 완료를 확인하고 넷째ㆍ다섯째는 종료 상태에서 미해제 의무 및 선행조건 위반 사유를 확인한다. 진단 경로의 \texttt{first\_bad\_idx/first\_missing\_idx}, 관련 의무 인덱스를 원 사건/의무 ID로 되돌린다. 조건 B는 A에 대응하는 관측 배열에 들어가지 않는다.

기존 모델은 P1--P4 플래그별 \texttt{A[] not flag}, \texttt{A[] not sticky\_unknown}, \texttt{A<> Observer.Done}을 사용했다. 최초 반례 접두 경로가 뒤에 생긴 미상 사유를 담지 못한 24건은 종료 상태 추가 질의로 진단했다. 원 파서는 변경하지 않았다. 두 기존 구현의 일치는 고정 규칙의 구현 일치성이지 그 규칙의 독립 타당성이나 자연어 해석의 정답성이 아니다.

보조 반응 모델은 기록 사건과 자차 속도 반응을 재생하고 제한 시간ㆍ반복 시계ㆍ출처 덮어쓰기를 검사한다. 중심 조건별 모델과 질문 및 입력이 다르므로 그 성공을 중심 문장 검사의 안전 검증으로 합치지 않는다.
\end{revision}
